\documentclass[10pt,a4paper]{amsart}
\usepackage[margin=19mm]{geometry}
\usepackage{amsmath,amssymb,mathtools}
\usepackage[scr=rsfs,cal=euler]{mathalfa}
\usepackage{xfrac}
\usepackage[unicode,hidelinks,bookmarks=false]{hyperref}
\newcommand{\ie}{i.e.\ }
\newcommand{\cf}{cf.\ }
\newcommand{\resp}{respectively\ }
\newcommand{\Subsection}{\subsection}
\providecommand{\nobreakdash}{\nobreak\hbox{-}}
\setlength{\emergencystretch}{2em}
\multlinegap0pt
\usepackage{float}
\usepackage{xcolor}
\newtheorem{theorem}{Theorem}
\title{Technical Note on Relating Scores of Tilted Distributions}
\author{Curtis McDonald}
\date{Source: arXiv:2604.27196v1 (CC BY 4.0)}
\begin{document}
\maketitle

\noindent\emph{Source and licence.} Technical Note on Relating Scores of Tilted Distributions. Version arXiv:2604.27196v1 (CC BY 4.0), distributed by arXiv under the Creative Commons Attribution 4.0 International licence, http://creativecommons.org/licenses/by/4.0/, as recorded in the arXiv licence metadata for this exact version and shown on the abstract page https://arxiv.org/abs/2604.27196v1. Full licence notice: \texttt{licenses/arxiv-2604.27196-CC-BY-4.0.txt}.\par\smallskip

\noindent\textit{Editorial scope.} Counted unit: the article's theorem main\_denoiser\_thm (source0.tex lines 121--139). The article's complete original proof of it is delivered with it (source0.tex lines 140--142, one \texttt{proof} environment of 3 lines). No mathematics is written, repaired or re-derived: the statement and the proof are the article's own lines, in the article's own notation, and no retained line is substituted or re-flowed.  No further source-local context is needed: every cross-reference of every retained line resolves inside the two delivered spans.  Nothing was omitted from any retained span, so the package carries no omission record.  No retained line contains a citation, so the wrapper carries no bibliography and no bibliography entry.  No retained line contains a figure environment, an image-inclusion command or drawing code, so no image is delivered and the package declares no asset dependency.  Editorial formatting only, in addition: an AMS \texttt{amsart} preamble that restates the article's own declarations and macros used by the retained lines, namely the environments \texttt{theorem} on separate counters, together with the standard packages those lines need.  The retained environments print in this document as Theorem 1 in order of appearance, whereas the article prints the same items with the numbering of its own full document.  The complete native source of the article as distributed in the arXiv e-print for this exact version is bundled unchanged as \texttt{sources/199454/source0.tex}.
\begin{theorem}\label{main_denoiser_thm}
Let $p(x)$ have denoiser $F[u,\sigma]$. Let $p^{*}(x)$ be a linear tilt by a vector $v$ to $p$ and a constant negative diagonal tilt with scaling $s \geq 0$,
\begin{align}
p^{*}(x)&= \frac{p(x)e^{v^{T}x - \frac{1}{2}s \|x\|^{2}}}{E_{p}[e^{v^{T}x - \frac{1}{2}s \|x\|^{2}}]} \label{const_neg_daig_tilt}
\end{align}
Assume $X' \sim p^{*}$ and for any $\sigma \in [0,1]$ define random variable
\begin{align}
U' = \sqrt{1-\sigma^{2}}X'+\sigma Z, Z \sim N(0,I)
\end{align}
Denote the denoiser for $p^{*}(x)$ as $G[u,\sigma]$. Then the denoiser for $p^{*}$ is the denoiser for $p$ evaluated at an affine mapping of the inputs. Define the time shift and location shift,
\begin{align}
&u'=\frac{\sigma^{2}}{\sqrt{(1-\sigma^{2}+s \sigma^{2})(1+s \sigma^{2})}} v+ \sqrt{\frac{1-\sigma^{2}}{(1-\sigma^{2}+s \sigma^{2})(1+s \sigma^{2})}} u \label{u_prime}\\ 
&\sigma'=\sqrt{\frac{\sigma^{2}}{1+s \sigma^{2}}} \label{sigma_prime}
\end{align}
Then the denoiser for $p^{*}$ is the denoiser for $p$ evaluated at these new values
\begin{align}
G[u,\sigma]= F[u', \sigma']
\end{align}
\end{theorem}

\begin{proof}
The proof is an algebra procedure in completing the square and collecting like terms. All terms involved are either linear or quadratic and combine under simple rules. See Appendix.
\end{proof}

\end{document}
